% 2010 年安徽卷文科第 9 题：正视图、侧视图与俯视图。
% 来源 PDF：普通高考/2010/2010安徽文.pdf，SHA-256 1bd928e09356787e648fa9efb9b4c1525f26334f65b08ba82fc736ebb110a4cb。
% 原页：第 1 页（A3 三栏）第 2 栏上部，书页码 571。
% 重绘依据：正视图总宽 8（上部宽 6、两侧各 1），上部高 8、底座高 2；
%   侧视图总宽 10（6、2、2），上部宽 2；俯视图外框 8×10，内框 6×2 位于距上边 6 处。
\begin{tikzpicture}[semithick, scale=0.72, every node/.style={inner sep=1.2pt, font=\small}]
  % 正视图
  \draw (0,0) rectangle (8,2);
  \draw (1,2) rectangle (7,10);
  \draw[<->] (0,10.5) -- (1,10.5) node[midway,above] {$1$};
  \draw[<->] (1,10.5) -- (7,10.5) node[midway,above] {$6$};
  \draw[<->] (7,10.5) -- (8,10.5) node[midway,above] {$1$};
  \draw[<->] (-0.45,2) -- (-0.45,10) node[midway,left] {$8$};
  \draw[<->] (-0.45,0) -- (-0.45,2) node[midway,left] {$2$};
  \node[below] at (4,-0.55) {正（主）视图};
  % 侧视图
  \draw (10,0) rectangle (20,2);
  \draw (16,2) rectangle (18,10);
  \draw[<->] (10,10.5) -- (16,10.5) node[midway,above] {$6$};
  \draw[<->] (16,10.5) -- (18,10.5) node[midway,above] {$2$};
  \draw[<->] (18,10.5) -- (20,10.5) node[midway,above] {$2$};
  \node[below] at (15,-0.55) {侧（左）视图};
  % 俯视图
  \draw (0,-14) rectangle (8,-4);
  \draw (1,-12) rectangle (7,-10);
  \draw[<->] (-0.45,-4) -- (-0.45,-10) node[midway,left] {$6$};
  \draw[<->] (-0.45,-10) -- (-0.45,-12) node[midway,left] {$2$};
  \draw[<->] (-0.45,-12) -- (-0.45,-14) node[midway,left] {$2$};
  \node[below] at (4,-14.55) {俯视图};
\end{tikzpicture}
